\section{Alexandrov Topology on Minkowski Space}\label{sctn:alexandrov-topology-on-minkowski-space}
To consider what that might mean let $p$ and $q$ be any two points in what physicists think of as Minkowski space. With these two points, one can define two natural sets: $I^+(p)$, the \textit{chronological future} of $p$ (i.e. the set of points reachable from $p$ via a future-directed, time-like curve), and $I^-(q)$, the \textit{chronological past} of $q$ (i.e. the set of all points that reach $q$ via a future-directed, time-like curve).

It turns out that $I^+(p)$ is an open set in the Euclidean topology for any $p$ in what physicists think of as Minkowski space (Proposition 2.8 of \href{https://doi.org/10.1137/1.9781611970609}{Penrose}). Similarly, the time-reversed version of Proposition 2.8 of \href{https://doi.org/10.1137/1.9781611970609}{Penrose} implies that $I^-(q)$ is an open set in that same topology. Hence, $I^+(p) \cap I^-(q)$ is also open in that topology.

This suggests that sets of the form $I^+(p) \cap I^-(q)$ might be able to be used as the \href{https://ncatlab.org/nlab/show/topological+base}{basis} of a new topology on $\mathbb{R}^4$. As it so happens, sets of this form \textit{can} be used as the basis for a topology on $\mathbb{R}^4$.

Sets $I^+(p) \cap I^-(q)$ with $p,q \in \mathbb{R}^4$ form the basis for what's known as the \textit{Alexandrov topology} of $\mathbb{R}^4$ (Definition 4.22 of \href{https://doi.org/10.1137/1.9781611970609}{Penrose}). (This is also sometimes known as the \textit{interval topology} of $\mathbb{R}^4$ (Remark 4.23 of \href{https://doi.org/10.1137/1.9781611970609}{Penrose}).) Furthermore, it turns out that the Alexandrov topology of $\mathbb{R}^4$ agrees with the Euclidean topology of $\mathbb{R}^4$ (Paragraph 4.23 of \href{https://doi.org/10.1137/1.9781611970609}{Penrose}), a welcome happenstance that ``explains'' why physicists had not confronted this earlier.

All that being said, we will equip $\mathbb{R}^4$ with the Alexandrov topology as it is the ``natural'' topology of relativity.  So for us \textit{Minkowski space} is $\mathbb{R}^4$ with the Alexandrov topology, the Minkowski metric, and standard time-orientation.

In this case the Euclidean topology and the Alexandrov topology agree, but on other manifolds this need not be the case.
